\section{Ausgangslage}

Durch zunehmende Digitalisierung ändern sich viele Prozesse in allen möglichen Branchen.
Besonders die Gastronomie ist davon sehr stark betroffen. Die Weise wie Dienstverhältnisse
in der Gastronomie eingegangen werden verändert sich stark. Vor allem Gastarbeiter und saisonale 
Dienstverhältnisse werden immer flexibler eingesetzt. Auch die Arbeitnehmer erwarten immer 
öfter Flexibilität.

Die suche nach Personal bei Engpässen durch Krankheit, Urlaub oder saisonale Wechsel
erfolgt oft über informelle Kanäle, telefonisch oder über unspezialisierte
Online-Plattformen. Diese Methoden führen oft zu Zeitineffizienz, schlechter Transparenz
und hohem Organisationsaufwand für Arbeitgeber.

Auch für Arbeitnehmer ist es so schwierig, kurzfristige Arbeitsmöglichkeiten zu finden.
Temporäre Stellenangebote sind nicht zentral und übersichtlich auffindbar, was Bewerbungen
erschwert und die Jobsuche zeitintensiv macht. Darüber hinaus ist es schwer Jobs zu
vergleichen, da kein einheitliches Format für Stellenangebote existiert.

Um diesen Problemen entgegenzuwirken, entwickelt das Unternehmen
\textit{Recode} eine App die es einfacher für Arbeitgeber macht, Arbeitskräfte während 
Personalengpässen zu finden und Arbeitnehmern erleichtert, spontane Dienststellen zu finden. Die App soll eine zentrale Plattform bieten, auf der Arbeitgeber kurzfristige Arbeitsmöglichkeiten posten und Arbeitnehmer diese schnell und einfach finden können. Durch die Digitalisierung dieses Prozesses sollen Effizienz, Transparenz und Benutzerfreundlichkeit verbessert werden.
Die App bringt Gastronomiebetriebe und Arbeitskräfte zusammen indem sie Stellenausschreibungen
und Bewerbungen auf einer Plattform bündelt. Sie sorgt für ein größeres Angebot für 
Arbeitnehmer und eine größere Reichweite für Arbeitgeber. Durch die einheitliche
Digitalisierung der Stellenausschreibungen können Jobs besser verglichen werden und der 
Bewerbungsprozess wird vereinfacht.

\section{Motivation und Problemstellung}

\subsection{Nutzergewinnung}

Um das beschriebene Problem der Gastronomie zu verbessern wird jedoch nicht nur die 
Technologie in Form der App Jobbing benötigt, sondern auch eine möglichst große Nutzerbasis.
Wenn die Plattform nur wenige Stellenangebote hat, sinkt auch die Attraktivität für
Arbeitnehmer, da die Auswahl zu gering ist, um sich vollständig auf die Vermittlungsplattform
zu verlassen. Auf der anderen Seite ist es auch von großer Bedeutung, eine große Anzahl an
Arbeitnehmern für die App zu gewinnen, da viele Arbeitskräfte benötigt werden,
um eine sichere Vermittlung während Personalengpässen und saisonalen Kundenüberlastungen zu
garantieren. Diese Sicherheit ist ausschlaggebend um informelle Kanäle zu ersetzen.
Deshalb ist es wichtig viele Nutzer für die App zu gewinnen, um eine kritische Masse zu
erreichen, die es ermöglicht, dass Arbeitgeber und Arbeitnehmer gleichermaßen von der
Plattform profitieren können.

Die Gewinnung dieser kritischen Masse ist eine zentrale Herausforderung. Die Strategie der
Informationsübermittlung, an die Kunden, ist dabei ein wichtiger Faktor, da sie in den 
frühen Phasen entscheidet ob die Nutzer der App vertrauen.

\subsection{Unsicherheit über optimale Informationsstrategie}

Das primäre Ziel der Werbestrategie ist es, die Bekanntheit der Plattform zu steigern und
neue Nutzer zu gewinnen. Dafür ist es wichtig, den potenziellen Nutzern möglichst viele
Informationen zu übermitteln, damit sie die Vorteile der Plattform verstehen.

In der Regel dienen Websites als erster Kontaktpunkt zwischen der Anwendung und den
Nutzern. Die Gestaltung der Website beeinflusst maßgeblich den ersten Eindruck, den sich 
die Kunden von der App machen. Daraus folgt auch die Entscheidung ob der/die Nutzer*in bereit ist
sich auf einen Download einzulassen, oder ob er/sie die Website ohne weitere Schritte verlässt.
Websites bieten viele Möglichkeiten um den User anzusprechen, es ist jedoch schwer sich von
abzuheben, da Endverbraucher täglich mit einer Vielzahl von Websites konfrontiert werden.

\subsection{Einfluss von Personalisierung und Tonalität}

Im Umgang mit der Gestaltung digitaler Dialogsysteme sind vor allem zwei Kriterien wichtig: die
Personalisierung und der Sprachstil. Die Personalisierung kann etwa durch das Einsetzen des Namens
oder durch eine individuelle Ansprache erfolgen. Diese Maßnahmen zielen darauf ab, die wahrgenommene
Relevanz zu steigern und eine intensivere Bindung zu schaffen.
Zugleich ist eine Personalisierung aber auch leicht als aufdringlich oder gar künstlich
empfindbar. Das gilt ebenso für die Tonalität eines Systems. Eine freundliche, um nicht zu sagen
„kumpelhafte“ Ansprache kann Nähe und Sympathie schaffen, während eine formelle, steife
Ansprache Professionalität und Seriosität signalisiert.
Welche dieser Varianten nun im konkreten Fall der Bewerbung einer Vermittlungsplattform für die
Gastronomie positivere Rückwirkungen entfaltet, ist nicht ohne weiteres abzuschätzen und kann nur
durch empirische Untersuchungen ermittelt werden.

\subsection{Fehlende empirische Grundlage}

Meistens werden Marketing- und UX-Entscheidungen nach dem Prinzip „trial and error“ oder
nach Best Practices getroffen, ohne die Alternativen systematisch und datengestützt zu
testen. Gerade im Projektkontext können verschiedene Varianten experimentell miteinander
verglichen und deren Wirkung auf das Nutzerverhalten und die Zufriedenheit beobachtet
werden. Hieraus kann dann eine Entscheidung abgeleitet werden, die über reine Vermutungen
hinausgeht.

\subsection{Forschungsleitende Fragestellungen}

Aus den dargestellten Überlegungen ergeben sich folgende
zentrale Fragestellungen:

\begin{itemize}
    \item Welche Form der Nutzerführung führt zu einer höheren
    Download-Bereitschaft der App?
    \item Wie wirkt sich Personalisierung im Chatbot-Dialog auf
    die wahrgenommene Relevanz aus?
    \item Beeinflusst die Tonalität des Chatbots die
    Zufriedenheit der Nutzer:innen?
    \item Welche Fragen und Bedenken äußern potenzielle
    Nutzer:innen vor einem Download?
    \item Ermitteln, welche Art von Funnel am wirksamsten ist.
    \item Untersuchen, welche Elemente die User-Experience positiv beeinflussen und die Zielgruppe überzeugen.
    \item Analyse des gezielten UI-Einsatzes, um Nutzer:innen effizient durch den Funnel zu führen.
    \item Evaluierung, welches Tracking-Tool für die Analyse des Nutzerverhaltens am besten geeignet ist.
\end{itemize}

\section{Zielsetzung}

Ziel dieser Diplomarbeit ist die Entwicklung und Konzeption zweier verschiedener
Informationsstrategien für die digitale App Vermarktung der App \textit{Jobbing}.
Dabei werden ein Chatbot als Dialogsystem und ein Multi-Step-Funnel implementiert und
miteinander verglichen.
Beide Strategien verfolgen das Ziel, die User von der App zu überzeugen und im Endeffekt
zu einem Download zu bewegen. Beide Strategien sollen dabei von Endverbrauchern getestet 
werden. Dabei werden Daten über das Nutzerverhalten erhoben, welche später ausgewertet werden,
um die verschiedenen Konzepte sinnvoll miteinander vergleichen zu können.

Teil der Arbeit ist auch die Ermittlung des optimalen Weges der Nutzerdatenerhebung und Auswertung.
Wichtig für den Vergleich der Systeme ist nicht nur ob der Nutzer einen Download getätigt
hat, sondern auch Nutzerverhalten das sich vor dem Download ereignet. In dieser Arbeit wird
evaluiert welche der Daten am aussagekräftigsten für den Vergleich sind und auch wie sie 
miteinander verglichen werden.

Außerdem soll festgestellt werden wie die Nutzerdaten am besten erhoben werden können. 
Dafür gibt es viele verschiedene Möglichkeiten. Dazu gehören eigene Implementierungen 
die genau auf den Einsatzzweck zugeschnitten werden können, oder bereits bewährte 
Analyseapplikationen die Entwicklern Von Unternehmen wie Google zur Verfügung gestellt
werden. Um die Daten anschaulich vergleichen zu können, kann auch eine grafische Darstellung
von Vorteil sein. Diese kann ebenfalls selbst implementiert werden, es können aber auch
Grafiken von Beispielsweise Google Analytics zur Evaluierung verwendet werden.

Um den Vergleich der zwei Varianten möglichst aussagekräftig zu machen, müssen die einzelnen
Werbekonzepte selbst möglichst effektiv konzipiert und implementiert werden. Dafür müssen 
auch die Einzelnen Varianten getestet und optimiert werden. Deshalb wird nochmals eine 
spezifische Zielsetzung definiert.

\subsection{Zielsetzung Chatbot}

Das Dialogsystem soll primär präzise Auskunft über die App geben und eine 
natürliche Konversation ermöglichen. Es soll dem Nutzer das Bild einer kompetenten
Fachkraft vermitteln, die hilfsbereit sinnvolle Antworten gibt. Darüber hinaus
soll der Assistent auch versuchen beim Nutzer Interesse zu wecken, und anhand
der Konversation die Vorteile von Jobbing zu vermitteln. Ein weiteres Ziel ist
es den User zum Download der App zu bewegen. Der Bot soll die Frage nach einem Download
gezielt erkennen und dann dem Nutzer die Möglichkeit zum Download geben. Diese Intention
soll erkannt werden, um später eine sinnvolle Auswertung zu ermöglichen.

Außerdem soll der Chatbot so implementiert werden, dass die Forschungsfragen
bezüglich Personalisierung und Tonalität beantwortet werden können. Dafür soll
das Dialogsystem so gestaltet werden, dass es auf verschiedene Arten 
kommunizieren kann, und diese getestet werden können. Für die Evaluierung 
dieser Fragen, ist es wichtig, dass die Nutzerdaten unverfälscht erhoben werden
können.

Die Zielsetzung besteht somit darin, eine datenbasierte
Bewertung des Chatbots zu ermöglichen und
Gestaltungsempfehlungen für dialogische
Marketing-Strategien abzuleiten.

\subsection{Zielsetzung Multi-Step-Funnel}

\section{Vorgehensweise}